% !PRINTED 519
\weberpage{519}
%<<CONT
holds, provided we set the exponent equal to zero when one of the prime factors does not divide the functional in question. A common divisor of the numbers $\alpha, \beta, \ldots$ can contain no prime factors other than $\pi_1, \pi_2, \ldots$, and every common divisor of $\alpha, \beta, \ldots$ has the form
\weq{8}{\varepsilon \delta = \pi_1^a \pi_2^b \ldots,}
in which $a$ may not be greater than the least of the numbers $a_1, b_1, \ldots$, and $b$ not greater than the least of the numbers $a_2, b_2, \ldots$, and so forth.

If $a$ is the least among the numbers $a_1, b_1, \ldots$, and $b$ the least among the numbers $a_2, b_2, \ldots$, then the number $\delta$ represented in (8) is the greatest common divisor of the numbers $\alpha, \beta, \ldots$ Expressed in words:

One obtains the greatest common divisor of several integral functionals $\alpha, \beta, \ldots$ by forming a product of prime factors, into which one takes up each prime factor as often as it divides each of the numbers $\alpha, \beta, \ldots$.

Two numbers are relatively prime when they contain no common prime factor.

Correspondingly we define as the \gesperrt{least common multiplum} $\mu$ of the functionals $\alpha, \beta, \ldots$ a product of prime factors, into which we take up a factor $\pi$ so often that it divides none of the functionals $\alpha, \beta, \ldots$ more often than it divides $\mu$. This functional $\mu$ then has, just like every functional associated with $\mu$, the twofold and, as we may add, characteristic property that it is divisible by each of the functionals $\alpha, \beta, \ldots$, and that every other functional which is divisible by $\alpha, \beta, \ldots$ is also divisible by $\mu$.

The least common multiplum of two relative prime functionals is their product.

These theorems can be applied in order to represent fractional functionals \gesperrt{in the simplest shape, or as reduced fractions}, by decomposing numerator and denominator into their prime factors and cancelling the greatest common divisor. Numerator and denominator of a reduced fraction are completely determined by the fraction itself, apart from a common unit factor, which remains undetermined.

In the same way one can bring any number of given fractions to a common denominator, the principal denominator, by choosing a common multiplum of all the given denominators as the common denominator.
